\section{Conclusion}

This study presents PRISM, a data platform that extracts plastic waste-related records from heterogeneous sources using large language models and other tools. It also introduces AUDA, an analytical platform that migrates validated records from PRISM into structured analytical tables and produces reproducible modeling workflows. We used these workflows to evaluate correlation and feature importance, missing value reconstruction, univariate temporal forecasting, pooled neural network forecasting, and country-specific multivariate forecasting.

The experiment results do not point to one model that works best in every situation. Support Vector Regression (SVR) performs well in global plastic production forecasting, which is a single-source time-series dataset with over 70 samples; it also performs relatively well in both reconstruction tasks. However, simple interpolation, smoothing, and statistical baselines remain competitive for short country-level series, where machine learning models can easily overfit or become unstable. Multivariate experiments show that pooled neural network pretraining can assist forecasting to some extent, but it does not show a qualitative improvement in multivariate forecasting tasks over traditional machine learning models such as Gaussian process regression. Since most country-level time-series datasets used in this study are short and noisy, the experimental results are provided for reference only.

The proposed selection efficiency coefficient for the one standard error rule provides a practical way to balance validation performance alongside model simplicity. In the reconstruction sensitivity study, when the selection efficiency is set to \(c_{\mathrm{SE}} = 0.1\), the selected hyperparameter configuration is often different from the one with the best validation performance, but the WAPE penalty is very small. A larger selection efficiency expands the candidate range, resulting in the selection of simpler but underfitting configurations. Overall, the results suggest that a modest selection efficiency coefficient can reduce the probability of choosing an overfitting configuration without hurting accuracy.

Data from different reports may be generated using different approaches, which can lead to inconsistencies. Future work should consider the metadata associated with each migrated data point, such as the source document and organization, and propose a new architecture for measuring the relationship between each value and its metadata.
